\section*{Capítulo \ref{ch:b3:gfd-foundations} --- Dinâmica de fluidos geofísicos}

\begin{solution}{exo:b3:gfd-foundations:1}
(a)~$f = 2\Omega\sin\varphi$: \num{7.292e-5}, \num{1.031e-4} e
\qty{1.263e-4}{s^{-1}}; $2\pi/f$: 23.93, 16.92 e \qty{13.82}{h}.
(b)~$U/f$: 2.74, 1.94 e \qty{1.58}{km}. (c)~A \ang{30}, $\sin\varphi =
1/2$ e $2\pi/f = 2\pi/\Omega$: o período inercial é o dia sideral,
\qty{23.93}{h}.
\end{solution}

\begin{solution}{exo:b3:gfd-foundations:2}
Com $f = \qty{1.031e-4}{s^{-1}}$: redemoinho de poeira \num{9.7e3}, brisa marítima 2.4,
depressão 0.15, giro \num{2.4e-4}. A depressão e o giro estão próximos do
equilíbrio geostrófico; o redemoinho de poeira gira nos dois sentidos, alheio
à rotação da Terra, e a brisa marítima só a percebe após um dia.
\end{solution}

\begin{solution}{exo:b3:gfd-foundations:3}
$\partial u_g/\partial z = (g/fT)\,|\partial T/\partial y| = 9.81\times
\num{5e-6}/(\num{1.031e-4}\times250) = \qty{1.90e-3}{s^{-1}}$; em
\qty{10}{km}, o vento de oeste aumenta \qty{19}{m/s}.
\end{solution}

\begin{solution}{exo:b3:gfd-foundations:4}
$E = \num{50e-8}$: $N = \qty{2.21e-3}{s^{-1}}$, um período de \qty{47}{min},
\num{1.27} ciclos por hora. $E = \num{1000e-8}$: $N = \qty{9.90e-3}{s^{-1}}$,
\qty{10.6}{min}, 5.7 ciclos por hora. $E = \num{e-8}$: $N =
\qty{3.13e-4}{s^{-1}}$, \qty{5.6}{h}, 0.18 ciclo por hora.
\end{solution}

\begin{solution}{exo:b3:gfd-foundations:5}
(a)~$W \sim UH/L = \qty{0.1}{m/s}$ e $\mathrm Dw/\mathrm Dt \sim UW/L =
\qty{e-6}{m/s^2}$, \num{1.0e-7} de $g$. O termo de pressão, $fUL/H =
\qty{0.103}{m/s^2}$, é maior por \num{1e5}: a razão, \num{9.7e-6}, é
$\mathrm{Ro}(H/L)^2$ com $\mathrm{Ro} = 0.097$. (b)~$\tilde fu =
2\Omega\cos\ang{45}\times10 = \qty{1.03e-3}{m/s^2}$, \num{1.05e-4} de $g$.
\end{solution}

\begin{solution}{exo:b3:gfd-foundations:6}
$fr = \qty{30}{m/s}$. (a)~$\mathrm{Ro}_g = 0.667$, $V = 40/(1 + \sqrt{3.667})
= \qty{13.7}{m/s}$, um terço mais lento que o geostrófico. (b)~$\mathrm{Ro}_g =
0.167$, $V = 10/(1 + \sqrt{0.333}) = \qty{6.34}{m/s}$. (c)~$V_g \le fr/4 =
\qty{7.5}{m/s}$, e então $V = fr/2 = \qty{15}{m/s}$, o dobro da
velocidade geostrófica.
\end{solution}

\begin{solution}{exo:b3:gfd-foundations:7}
(a)~$\Omega = 2\pi/6 = \qty{1.047}{rad/s}$, $f = \qty{2.094}{s^{-1}}$.
$\mathrm{Ro} = 0.01/(2.094\times0.05) = 0.095$ e $\mathrm{Ek} =
\num{e-6}/(2.094\times0.04) = \num{1.19e-5}$: ambos pequenos, de modo que o
escoamento obedece ao \omterm{thm:b3:gfd-foundations:taylor-proudman}{teorema de Taylor--Proudman} e forma-se uma \omterm{thm:b3:gfd-foundations:taylor-proudman}{coluna de Taylor}. (b)~$0.2/\sqrt{\num{e-6}
\times1.047} = \qty{195}{s}$, \qty{3.3}{min}, 33 voltas da mesa.
\end{solution}

\begin{solution}{exo:b3:gfd-foundations:8}
$\partial\bar T/\partial y = -f\,\Delta u/(R\ln(850/250))$, com
$\ln(850/250) = 1.224$. A \ang{30} N, $-\num{7.29e-5}\times36.12/(286.9
\times1.224) = \qty{-7.50}{K}$ por \qty{1000}{km}; a \ang{45} N,
$-\num{1.031e-4}\times17.18/351.1 = \qty{-5.05}{K}$ por \qty{1000}{km}. O
cisalhamento é proporcional ao gradiente dividido por $f$: o gradiente é 1.5
vezes maior a \ang{30} e $f$ 1.4 vezes menor, de modo que o cisalhamento é 2.1 vezes
maior.
\end{solution}

\begin{solution}{exo:b3:gfd-foundations:9}
(a)~$L_R = 0.0113\times11\,000/\num{1.031e-4} = \qty{1206}{km}$;
$0.00357\times1000/\num{9.47e-5} = \qty{37.7}{km}$, 32 vezes menor.
(b)~$\beta = 2\Omega\cos\varphi/a$: \num{1.62e-11} e
\qty{1.74e-11}{m^{-1}.s^{-1}}. $\beta L_R/f = 0.19$ para a atmosfera, 0.007
para o oceano.
\end{solution}

\begin{solution}{exo:b3:gfd-foundations:10}
(a)~$\mathrm dw/\mathrm dt = -(r + \mathrm if)w$, logo $w = U\mathrm
e^{-rt}\mathrm e^{-\mathrm ift}$: a velocidade gira à taxa $f$ enquanto
sua magnitude decai como $\mathrm e^{-rt}$. A posição $x + \mathrm iy =
U(1 - \mathrm e^{-(r + \mathrm if)t})/(r + \mathrm if)$ espirala em direção ao
ponto $U/(r + \mathrm if)$. (b)~$x + \mathrm iy = U(1 - \mathrm e^{-(r +
\mathrm if)t})/(r + \mathrm if) + u_0t$: a espiral se desdobra em laços.
(c)~$0.2\,\mathrm e^{-1} = \qty{0.074}{m/s}$; o centro derivou $u_0t =
0.05\times4\times86\,400 = \qty{17.3}{km}$ para leste.
\end{solution}

\begin{solution}{exo:b3:gfd-foundations:11}
$f = \qty{9.37e-5}{s^{-1}}$; $\partial u/\partial z = (g/f\rho_0)\,\partial
\rho/\partial y = 9.81\times10^{-5}/(\num{9.37e-5}\times1025) =
\qty{1.02e-3}{s^{-1}}$. Com $u = 0$ a \qty{1000}{m}, a corrente superficial é
\qty{1.02}{m/s}, para leste: no Hemisfério Norte a corrente flui com a água
leve à sua direita.
\end{solution}

\begin{solution}{exo:b3:gfd-foundations:12}
(a)~$N^2 = (g/\rho)\,\Delta\sigma_\theta/\Delta z$. 0--\qty{100}{m}:
$9.81\times0.668/(1026.4\times100)$, $N = \qty{7.99e-3}{s^{-1}}$, 4.6
ciclos por hora; 100--\qty{500}{m}: \qty{3.02e-3}{s^{-1}}, 1.7 ciclos por
hora; 500--\qty{1000}{m}: \qty{3.24e-3}{s^{-1}}, 1.9 ciclos por hora.
(b)~Os \qty{100}{m} superiores: contêm a água superficial quente acima da
picnoclina sazonal. Entre 100 e \qty{500}{m} a água é mais
uniforme do que abaixo. (c)~$E = N^2/g$ é \num{650e-8}, \num{93e-8} e
\qty{107e-8}{m^{-1}}, todos dentro da faixa do oceano superior.
\end{solution}

\begin{solution}{pb:b3:gfd-foundations:1}
\textbf{1.} $f = 2\times\num{7.292115e-5}\sin\ang{48} =
\qty{1.084e-4}{s^{-1}}$; $2\pi/f = \qty{16.1}{h}$.

\textbf{2.} Banheira \num{3.1e3}; tanque ($f = 4\pi/10 = \qty{1.257}{s^{-1}}$)
0.020; brisa 2.3; depressão 0.14; giro \num{2.3e-4}.

\textbf{3.} A banheira não sente nada da rotação da Terra; a brisa
só a percebe após algumas horas. O tanque, a depressão e o giro têm
\omterm{def:b3:gfd-foundations:numbers}{números de Rossby} pequenos e estão próximos do equilíbrio geostrófico.

\textbf{4.} Tanque $\num{e-6}/(1.257\times0.04) = \num{2.0e-5}$; depressão
$10/(\num{1.084e-4}\times10^8) = \num{9.2e-4}$; giro
$\num{e-3}/(\num{1.084e-4}\times10^6) = \num{9.2e-6}$.

\textbf{5.} $\delta/H = \sqrt{2\nu/(fH^2)} = \sqrt{2\,\mathrm{Ek}}$. Tanque
$\sqrt{2\times\num{e-6}/1.257} = \qty{1.3}{mm}$; ar $\sqrt{20/\num{1.084e-4}}
= \qty{430}{m}$.

\textbf{6.} O atrito atua em camadas de alguns milésimos da profundidade no
tanque e no oceano, e de alguns centésimos na atmosfera; acima delas o
escoamento é quase sem atrito.

\textbf{7.} $\rho = p/(RT) = 101\,290/(286.9\times288.1) =
\qty{1.225}{kg/m^3}$.

\textbf{8.} $u_g = (1/\rho f)\,|\partial p/\partial y| = (400/300\,000)/(1.225
\times\num{1.084e-4}) = \qty{10.0}{m/s}$, de oeste, com a baixa
pressão ao norte à esquerda.

\textbf{9.} $\mathrm{Ro}_g = 10.04/(\num{1.084e-4}\times600\,000) = 0.154$;
$V = 2\times10.04/(1 + \sqrt{1.618}) = \qty{8.84}{m/s}$, \qty{12}{\percent}
mais lento.

\textbf{10.} $V = 2\times10.04/(1 + \sqrt{0.382}) = \qty{12.4}{m/s}$. O
maior $V_g$ é $fr/4 = \qty{16.3}{m/s}$, um gradiente $\rho fV_g =
\qty{2.16e-3}{Pa/m}$, \qty{6.5}{hPa} a cada \qty{300}{km}.

\textbf{11.} $g/c_p = \qty{9.77}{K/km}$. $N^2 = (9.81/255.7)\times(9.77 -
6.49)\times10^{-3} = \qty{1.26e-4}{s^{-2}}$, $N = \qty{0.0112}{s^{-1}}$, com
período de \qty{9.3}{min}.

\textbf{12.} Com $\Gamma = 0$, $N^2 = (g/T)(g/c_p) = g^2/(c_pT) =
\qty{4.43e-4}{s^{-2}}$, $N = \qty{0.0210}{s^{-1}}$, 1.9 vezes o
valor troposférico.

\textbf{13.} $L_R = 0.0112\times10^4/\num{1.084e-4} = \qty{1035}{km}$, o
tamanho da depressão.

\textbf{14.} $g' = \qty{0.0147}{m/s^2}$; $\sqrt{g'h} = \qty{0.54}{m/s}$;
$L_R = 0.542/\num{1.084e-4} = \qty{5.0}{km}$.

\textbf{15.} Sim: o lago, com \qty{12}{km} de largura, tem mais que o dobro
de seu raio de Rossby, de modo que suas ondas internas e correntes são
desviadas pela força de Coriolis.

\textbf{16.} $(249.4 - 253.0)/600 = \qty{-6.0}{K}$ a cada \qty{1000}{km}.

\textbf{17.} $\ln(1013/265) = 1.3409$. Sul:
$286.9\times253.0\times1.3409/9.81 = \qty{9922}{m}$; norte \qty{9781}{m}.
A superfície de \qty{265}{hPa} fica \qty{141}{m} mais baixa no norte, uma
inclinação de \num{2.35e-4} em relação à superfície inferior.

\textbf{18.} $\Delta u = (g/f)\times\text{inclinação} = 9.81\times\num{2.35e-4}/
\num{1.084e-4} = \qty{21.3}{m/s}$.

\textbf{19.} $(R/f)\ln(p_b/p_t)\,|\partial\bar T/\partial y| =
(286.9/\num{1.084e-4})\times1.3409\times\num{6.0e-6} = \qty{21.3}{m/s}$, o
mesmo resultado: os dois cálculos são um só.

\textbf{20.} \omterm{def:b3:gfd-foundations:baroclinic}{Baroclínico}: a temperatura, e portanto a densidade, muda ao
longo das superfícies de pressão. \omterm{def:b3:gfd-foundations:baroclinic}{Barotrópico}, o vento não mudaria com a
altura: \qty{10.0}{m/s}.

\textbf{21.} $(g/f\bar T)\,|\partial T/\partial y|\,\Delta z = 9.81\times
\num{6.0e-6}\times10^4/(\num{1.084e-4}\times255.7) = \qty{21.2}{m/s}$,
dentro de \qty{0.3}{\percent}.

\textbf{22.} $10.04 + (286.9/\num{1.084e-4})\times\num{6.0e-6}\times
\ln(1013/541) = 10.04 + 9.96 = \qty{20.0}{m/s}$.

\textbf{23.} $\partial u/\partial z = 21.3/10^4 = \qty{2.13e-3}{s^{-1}}$;
$\mathrm{Ri} = \num{1.26e-4}/\num{4.54e-6} = 28$, muito acima de 0.25: o
cisalhamento não torna a troposfera livre turbulenta.

\textbf{24.} O vento de Orrin, \qty{31}{m/s}, é mais forte que a média
zonal a \ang{45} N e mais fraco que a \ang{30} N. A corrente de jato da frente polar
oscila para norte e para sul; em cada instante seu gradiente e seu vento se
concentram numa faixa, e a média ao longo de um círculo de latitude os espalha.

\textbf{25.} $10.0 + 21.3$: o vento geostrófico a \qty{10}{km} acima da
estepe é \textbf{\qty[text-series-to-math=true]{31}{m/s}}, de oeste
(\qty{31.3}{m/s} pelo cálculo).
\end{solution}
